\subsection*{الف}
می‌دانیم اگر
$y[n] \longleftrightarrow Y(e^{j\omega})$
آنگاه با استفاده از خاصیت مزدوج مختلط و وارونگی زمان خواهیم داشت
\begin{flalign*}
y^*[-n] \longleftrightarrow Y^*(e^{j\omega}). &&
\end{flalign*}
حال با قضیهٔ کانولوشن در \lr{DTFT} اگر
$Z(e^{j\omega}) = X(e^{j\omega})\,Y^*(e^{j\omega})$
باشد، آنگاه
\begin{flalign*}
z[n] = x[n] * y^*[-n]
= \sum_{k=-\infty}^{\infty} x[k]\,y^*[k-n]. &&
\end{flalign*}

\subsection*{ب}
از رابطهٔ معکوس \lr{DTFT} برای
$z[n]$
داریم
\begin{flalign*}
z[n] = \frac{1}{2\pi}\int_{-\pi}^{\pi} Z(e^{j\omega})e^{j\omega n}\,d\omega
= \frac{1}{2\pi}\int_{-\pi}^{\pi} X(e^{j\omega})Y^*(e^{j\omega})e^{j\omega n}\,d\omega. &&
\end{flalign*}
با قرار دادن $n=0$:
\begin{flalign*}
z[0] = \frac{1}{2\pi}\int_{-\pi}^{\pi} X(e^{j\omega})Y^*(e^{j\omega})\,d\omega. &&
\end{flalign*}
از قسمت (الف) نیز
\begin{flalign*}
z[0] = \sum_{k=-\infty}^{\infty} x[k]\,y^*[k]. &&
\end{flalign*}
پس نتیجه می‌شود
\begin{flalign*}
\sum_{n=-\infty}^{\infty} x[n]y^*[n]
= \frac{1}{2\pi}\int_{-\pi}^{\pi} X(e^{j\omega})Y^*(e^{j\omega})\,d\omega. &&
\end{flalign*}